\documentclass[12pt]{article}

\usepackage[T1]{fontenc}
\usepackage[utf8]{inputenc}
\usepackage{lmodern}
\usepackage[a4paper, margin=3cm]{geometry}

\title{Coordinating Multilayer Power System Flexibility:\\
A Hybrid MILP--GNN--EBM Framework for Large-Scale Scenario Exploration}
\author{Th\'eotime Coudray* \\ St\'ephane Goutte*\\[2pt]
\normalsize *Universit\'e de Versailles Saint-Quentin-en-Yvelines (UVSQ), UMI SOURCE}
\date{}

\begin{document}

\maketitle

\begin{abstract}
The large-scale integration of variable renewable energy, storage, and demand response
requires coordinating flexibility across multiple spatial layers and temporal scales.
Operators and planners increasingly need to explore thousands of what-if scenarios---spanning
extreme weather, topology changes, and evolving flexibility portfolios---rather than relying on
single-shot optimisation. Mixed-Integer Linear Programming (MILP) delivers rigorous and auditable
solutions to such problems, but repeated solves become computationally prohibitive as scenario
volume and model complexity grow. This paper proposes a hybrid MILP--GNN--EBM methodology designed
for near-real-time scenario exploration. A synthetic scenario generator produces diverse
power-system instances, ranked by an
ex-ante criticality index combining physical stress and combinatorial hardness. A multi-layer MILP
oracle solves each instance to produce optimal binary commitment and continuous dispatch labels.
Scenarios are then projected onto heterogeneous temporal graphs and encoded by a self-supervised
Hierarchical Temporal Encoder into compact topology-aware embeddings. Conditioned on these
embeddings, an Energy-Based Model learns an energy landscape over binary activation decisions
and generates diverse binary candidates. A feasibility decoder repairs residual constraint violations, and a two-stage LP worker recovers
the optimal continuous dispatch, guaranteeing complete and economically meaningful operating plans.
Evaluation on 300 out-of-sample scenarios across three criticality families shows that the
pipeline's advantage scales positively with scenario difficulty. On high-criticality instances, the pipeline achieves a median cost gap of
$-0.04\%$ relative to the MILP, with a median speedup of $11.5\times$ on the 41 scenarios
where the MILP hits its time limit.
\end{abstract}

\end{document}
